\section{Related Work}
\label{sec:related}
\textbf{Exploration Based on Novelty Bonuses} Learning from sparse or otherwise difficult-to-optimize reward functions is a long-standing problem in reinforcement learning.
%
There is a large body of work which defines intrinsic rewards based on novelty bonuses \citep{schmidhuber1991possibility, E3, Rmax, stadie2015incentivizing,  PseudoCounts, ICM, RND, max, RIDE, ecoffet2019go, PCPG, NovelD, E3B,lu2024intelligent}. These methods tend to make minimal assumptions about the task at hand, operate online without requiring external data, and sometimes come with theoretical guarantees. However, since they are fundamentally \textit{tabula-rasa}, they must rediscover much of the structure in the task that might already be encoded as prior knowledge in an LLM. Therefore, they tend to have difficulty exploring environments of very high complexity, such as NetHack, in a tractable amount of time \citep{klissarovdoro2023motif}.
\looseness=-1

\textbf{LLM-aided Reward Design}
In addition to Motif~\citep{klissarovdoro2023motif}, several works have sought to leverage the prior knowledge encoded in LLMs to produce intrinsic rewards.
Eureka \citep{ma2023eureka}, Auto-MC \citep{li2024auto}, L2R~\citep{yu2023language} and Text2Reward~\citep{xie2023text2reward}
all use LLMs to generate executable code which computes intrinsic rewards from the underlying environment state, conditioned on a task description. The generated reward function code is then iteratively improved based on aggregate statistics from agents trained with the current reward.
%
However, a disadvantage with intrinsic rewards represented as code is that they require access to an interpretable underlying state representation, and it is unclear how to leverage non-numerical features such as those provided by unstructured language captions.
The works of \citet{kwon2023reward, chu2023accelerating} also successfully used LLMs conditioned on task descriptions to directly generate binary rewards in an online manner, and did not train a reward model. This was possible due to evaluating on environments and tasks which could be solved with a relatively small number of observations, whereas we consider complex open-ended environments with billions of observations so that labeling them all with an LLM would be computationally infeasible. Also of note is the work of \citet{wu2024read}, which additionally conditioned LLMs on user manuals to define the intrinsic reward. \looseness=-1



\textbf{Goal-conditioned Reward Design}
A different approach to reward function design is to define rewards as the distance between the agent's current state and  the goal. For example, one line of work learns a state embedding in a self-supervised manner which converts geodesic distances in the original space to Euclidean distances in feature space  \citep{wu2018the, DBLP:journals/corr/abs-2107-05545, gomez2024properlaplacianrepresentationlearning}.
Another line of work of \citep{fan2022minedojo, rocamonde2023vision, adeniji2023languagerewardmodulationpretraining,kim2024guide} leverages pretrained image and text encoders, and defines rewards to be some measure of similarity between embeddings of visual observations and embeddings of textual task descriptions.
%
Using a question generation and answering system, \citet{carta2022eager} extract auxiliary objectives from the goal description and construct intrinsic rewards.
%
An interesting combination of goal-conditioned and LLM-aided reward design is the \ellm approach introduced in \citet{pmlr-v202-du23f}, which generates candidate goals and uses the distance in LLM embedding space to define the reward. 
%
%
This approach shares the limitations of the works of \citet{kwon2023reward, chu2023accelerating} discussed in the previous section, in that it requires an LLM call for each agent observation, which becomes computationally infeasible in high-throughput settings involving billions of observations.
%
We discuss more works that utilize LLM for RL broadly in \Cref{app:additional_related_work}.




%
%
%
%
